\chapter{Tầng 1: Thu thập và Chuẩn hóa Dữ liệu}

\section{Tổng quan Tầng Thu thập Dữ liệu}

Tầng 1 là lớp nền tảng của hệ thống STWI, có nhiệm vụ tiếp nhận dữ liệu thô từ phần cứng ngoại vi và chuyển hóa thành Tensor chuẩn hóa $\mathbf{X} \in \mathbb{R}^{B \times 12 \times 14}$ để cung cấp cho Tầng 2. Dữ liệu đến từ hai nguồn chính: (1) Mạng lưới camera CCTV và (2) Hệ thống cảm biến môi trường và khí tượng.

\section{Phân hệ Thị giác Máy tính (CCTV Pipeline)}

\subsection{Kiến trúc Pipeline CCTV}

\begin{figure}[H]
\centering
\begin{tikzpicture}[node distance=0.6cm and 0.5cm, scale=0.9, transform shape]
  \node[sysbox=tier1col, text width=2.8cm, minimum height=1.5cm] (cam)
    {Camera CCTV\\1.000 điểm\\25 FPS};
  \node[sysbox=tier2col, right=1.2cm of cam, text width=3cm, minimum height=1.5cm] (yolo)
    {YOLOv8\\Object Detection\\NVIDIA TensorRT};
  \node[sysbox=tier2col, right=1.2cm of yolo, text width=3cm, minimum height=1.5cm] (bt)
    {ByteTrack\\Multi-Object\\Tracking};
  \node[sysbox=tier3col, right=1.2cm of bt, text width=3.5cm, minimum height=1.5cm] (agg)
    {Aggregation\\Chu kỳ 5 phút\\/camera};
  \node[sysbox=tier4col, right=1.2cm of agg, text width=2.8cm, minimum height=1.5cm] (tsdb)
    {TimeSeries DB\\InfluxDB\\/ TimescaleDB};

  \draw[arrow] (cam) -- (yolo);
  \draw[arrow] (yolo) -- (bt);
  \draw[arrow] (bt) -- (agg);
  \draw[arrow] (agg) -- (tsdb);

  % Output labels
  \node[below=0.6cm of agg, font=\tiny\itshape, text=codegray, text width=3.5cm, align=center]
    {3 biến:\\traffic\_volume\\avg\_velocity\\vehicle\_class};
  \node[below=0.6cm of tsdb, font=\tiny\itshape, text=codegray, text width=2.8cm, align=center]
    {JSON records\\timestamped};
\end{tikzpicture}
\caption{Luồng xử lý dữ liệu CCTV --- từ Camera đến TimeSeries DB}
\label{fig:cctv_pipeline}
\end{figure}

\subsection{YOLOv8 -- Phát hiện Phương tiện}

YOLOv8 \cite{jocher2023yolo} là kiến trúc phát hiện vật thể (object detection) thời gian thực tiên tiến, xử lý mỗi khung hình (frame) trong một lần forward pass duy nhất:

\begin{itemize}[leftmargin=1.5cm, itemsep=4pt]
  \item \textbf{Backbone:} CSPDarknet53 với C2f modules cho trích xuất đặc trưng đa tỉ lệ.
  \item \textbf{Neck:} PANet (Path Aggregation Network) để kết hợp đặc trưng từ nhiều tỉ lệ.
  \item \textbf{Head:} Decoupled head dự báo bounding box và class score độc lập.
  \item \textbf{Tốc độ inference:} $\sim$3ms/frame trên NVIDIA A100 với TensorRT FP16.
\end{itemize}

Các lớp phương tiện cần nhận diện:

\begin{table}[H]
\centering
\caption{Phân loại phương tiện trong STWI CCTV Pipeline}
\label{tab:vehicle_class}
\renewcommand{\arraystretch}{1.4}
\begin{tabular}{>{\centering}p{1.5cm} p{3.5cm} p{5cm} p{2cm}}
\toprule
\textbf{Lớp ID} & \textbf{Loại phương tiện} & \textbf{Ví dụ cụ thể} & \textbf{PCU$^*$} \\
\midrule
0 & Xe máy & Xe ga, xe số, xe đạp điện & 0.5 \\
\rowcolor{rowgray}
1 & Ô tô con & Sedan, SUV, MPV & 1.0 \\
2 & Xe tải & Tải nhỏ, tải trung, container & 2.0--3.5 \\
\rowcolor{rowgray}
3 & Xe buýt & Buýt thường, buýt nhanh BRT & 2.5 \\
\bottomrule
\multicolumn{4}{l}{\small $^*$PCU: Passenger Car Unit --- đơn vị quy đổi ra xe con tương đương}
\end{tabular}
\end{table}

\subsection{ByteTrack -- Theo dõi Đa đối tượng}

ByteTrack \cite{zhang2020bytetrack} giải quyết vấn đề đếm trùng phương tiện giữa các frame liên tiếp bằng cách theo dõi \textit{tất cả} bounding box (bao gồm cả các detection có confidence thấp):

\begin{enumerate}[leftmargin=1.5cm, itemsep=4pt]
  \item \textbf{High-confidence detections} ($\text{conf} \geq \theta_{\text{high}}$): Ghép với tracks hiện có theo IoU (Hungarian algorithm).
  \item \textbf{Low-confidence detections} ($\theta_{\text{low}} \leq \text{conf} < \theta_{\text{high}}$): Thử ghép với các tracks bị miss ở bước trên.
  \item \textbf{Kalman Filter}: Dự báo vị trí track trong frame tiếp theo để xử lý occlusion.
\end{enumerate}

Kết quả: Mỗi phương tiện được gán ID duy nhất, cho phép tính lưu lượng thực (không đếm trùng) và vận tốc trung bình từ quỹ đạo di chuyển.

\subsection{Các Biến Đầu ra từ CCTV}

Sau mỗi chu kỳ 5 phút, pipeline CCTV xuất 3 biến cho mỗi nút giao:

\begin{table}[H]
\centering
\caption{Biến đặc trưng từ CCTV Pipeline}
\label{tab:cctv_features}
\renewcommand{\arraystretch}{1.4}
\begin{tabular}{p{3.5cm} p{2.5cm} p{6cm}}
\toprule
\textbf{Tên biến} & \textbf{Đơn vị} & \textbf{Mô tả} \\
\midrule
\texttt{traffic\_volume} & xe/5 phút & Tổng số phương tiện đi qua nút giao (PCU-weighted) \\
\rowcolor{rowgray}
\texttt{avg\_velocity} & km/h & Tốc độ trung bình của tất cả phương tiện \\
\texttt{vehicle\_class\_ratio} & $[0,1]$ & Tỷ lệ ô tô trong tổng lưu lượng \\
\bottomrule
\end{tabular}
\end{table}

\section{Phân hệ Cảm biến Môi trường và Khí tượng}

\subsection{Mạng lưới Cảm biến}

Hệ thống tích hợp dữ liệu từ hai loại cảm biến:

\begin{table}[H]
\centering
\caption{Danh sách biến cảm biến môi trường và khí tượng}
\label{tab:sensors}
\renewcommand{\arraystretch}{1.4}
\begin{tabular}{p{2cm} p{3.2cm} p{2.5cm} p{4cm}}
\toprule
\textbf{Nhóm} & \textbf{Biến} & \textbf{Ký hiệu} & \textbf{Đơn vị / Dải giá trị} \\
\midrule
\multirow{5}{*}{Khí thải} & Carbon monoxide & CO & $\mu$g/m$^3$, 0--500 \\
  & Carbon dioxide & CO$_2$ & ppm, 300--5.000 \\
  & Nitrogen oxides & NO$_x$ & $\mu$g/m$^3$, 0--200 \\
  & Bụi mịn nhỏ & PM$_{2.5}$ & $\mu$g/m$^3$, 0--250 \\
  & Bụi mịn lớn & PM$_{10}$ & $\mu$g/m$^3$, 0--400 \\
\midrule
\rowcolor{rowgray}
\multirow{3}{*}{Khí tượng} & Nhiệt độ & $T$ & $^\circ$C, 15--42 \\
\rowcolor{rowgray}
  & Độ ẩm & $H$ & \%, 30--100 \\
\rowcolor{rowgray}
  & Tốc độ gió & $W$ & m/s, 0--20 \\
\bottomrule
\end{tabular}
\end{table}

\begin{tcolorbox}[tipbox]
\textbf{Logic kỹ thuật quan trọng:} Tốc độ gió $W$ có mối quan hệ phi tuyến với PM$_{2.5}$ và PM$_{10}$: gió mạnh ($W > 5$ m/s) phát tán bụi nhanh, nhưng gió yếu ($W < 1$ m/s) cho phép bụi tích tụ. Mô hình STGCN cần nắm bắt được mối quan hệ này để dự báo chất lượng không khí trong điều kiện ùn tắc kéo dài.
\end{tcolorbox}

\subsection{Giao thức Thu thập và Xử lý Mất dữ liệu}

Cảm biến giao tiếp qua MQTT với broker trung tâm. Chiến lược xử lý dữ liệu thiếu:

\begin{itemize}[leftmargin=1.5cm, itemsep=3pt]
  \item \textbf{Ngắt $< 15$ phút:} Linear interpolation giữa hai điểm đo lân cận.
  \item \textbf{Ngắt 15--60 phút:} Sử dụng giá trị trung bình theo giờ tương ứng từ lịch sử 30 ngày.
  \item \textbf{Ngắt $> 60$ phút:} Đánh dấu flag `\texttt{sensor\_unavailable}', loại nút này khỏi graph inference.
\end{itemize}

\section{Cấu trúc và Chuẩn hóa Dữ liệu}

\subsection{14 Biến Đặc trưng (Feature Engineering)}

Sau khi thu thập, dữ liệu được ghép thành vector đặc trưng 14 chiều cho mỗi nút giao tại mỗi bước thời gian 5 phút:

\begin{table}[H]
\centering
\caption{14 biến đặc trưng đầu vào cho mô hình STGCN+LSTM}
\label{tab:features}
\renewcommand{\arraystretch}{1.4}
\begin{tabular}{>{\centering}p{1.5cm} >{\centering}p{0.8cm} p{3.5cm} p{5cm}}
\toprule
\textbf{Nhóm} & \textbf{ID} & \textbf{Tên biến} & \textbf{Mô tả} \\
\midrule
\multirow{3}{*}{Giao thông} & F1 & \texttt{traffic\_volume} & Lưu lượng xe (PCU/5 min) \\
  & F2 & \texttt{avg\_velocity} & Vận tốc trung bình (km/h) \\
  & F3 & \texttt{vehicle\_class\_ratio} & Tỷ lệ ô tô $(0, 1)$ \\
\midrule
\rowcolor{rowgray}
\multirow{5}{*}{Khí thải} & F4 & \texttt{co\_level} & Nồng độ CO ($\mu$g/m$^3$) \\
\rowcolor{rowgray}
  & F5 & \texttt{co2\_level} & Nồng độ CO$_2$ (ppm) \\
\rowcolor{rowgray}
  & F6 & \texttt{nox\_level} & Nồng độ NO$_x$ ($\mu$g/m$^3$) \\
\rowcolor{rowgray}
  & F7 & \texttt{pm25\_level} & Bụi mịn PM$_{2.5}$ ($\mu$g/m$^3$) \\
\rowcolor{rowgray}
  & F8 & \texttt{pm10\_level} & Bụi mịn PM$_{10}$ ($\mu$g/m$^3$) \\
\midrule
\multirow{3}{*}{Khí tượng} & F9 & \texttt{temperature} & Nhiệt độ ($^\circ$C) \\
  & F10 & \texttt{humidity} & Độ ẩm (\%) \\
  & F11 & \texttt{wind\_speed} & Tốc độ gió (m/s) \\
\midrule
\rowcolor{rowgray}
\multirow{3}{*}{Phụ trợ} & F12 & \texttt{hour\_sin} & $\sin(2\pi \cdot \text{hour}/24)$ \\
\rowcolor{rowgray}
  & F13 & \texttt{day\_of\_week} & Mã hóa vòng ngày trong tuần \\
\rowcolor{rowgray}
  & F14 & \texttt{signal\_phase} & Pha đèn tín hiệu hiện tại $\in \{0,1,2\}$ \\
\bottomrule
\end{tabular}
\end{table}

\subsection{Chuẩn hóa MinMaxScaler}

Do các features có thang đo rất khác nhau (nhiệt độ 15--42$^\circ$C vs lưu lượng 0--2.000 PCU), toàn bộ phải chuẩn hóa về dải $(0, 1)$:

\begin{equation}
  x_{\text{norm}} = \frac{x - x_{\min}}{x_{\max} - x_{\min}}
  \label{eq:minmax}
\end{equation}

\begin{tcolorbox}[dangerbox]
\textbf{Yêu cầu bắt buộc (QC-4):} Tham số $x_{\min}$, $x_{\max}$ được fit \textbf{chỉ trên tập training}. Tuyệt đối không fit lại khi inference để tránh data leakage. Scaler state phải được serialize (pickle) và lưu cùng với model checkpoint.
\end{tcolorbox}

\subsection{Xây dựng 3D Tensor}

\begin{figure}[H]
\centering
\begin{tikzpicture}[scale=0.85]
  % Draw 3D cuboid
  \def\dx{5.5} \def\dy{2.5} \def\dz{3.5}
  \def\ox{0.8} \def\oy{0.6}

  % Front face
  \draw[fill=bklightblue!50, draw=bkblue!60] (0,0) -- (\dx,0) -- (\dx,\dy) -- (0,\dy) -- cycle;
  % Top face
  \draw[fill=tier1col!15, draw=bkblue!60]
    (0,\dy) -- (\ox,\dy+\oy) -- (\dx+\ox,\dy+\oy) -- (\dx,\dy) -- cycle;
  % Right face
  \draw[fill=tier2col!15, draw=bkblue!60]
    (\dx,0) -- (\dx+\ox,\oy) -- (\dx+\ox,\dy+\oy) -- (\dx,\dy) -- cycle;

  % Dimension labels
  \draw[<->, thick, bkblue!80] (-0.4,0) -- (-0.4,\dy)
    node[midway, left, font=\small\bfseries] {Time Steps\\$T=12$\\(60 phút)};
  \draw[<->, thick, bkblue!80] (0,-0.5) -- (\dx,-0.5)
    node[midway, below, font=\small\bfseries] {Features $F=14$};
  \draw[<->, thick, bkblue!80] (\dx+\ox+0.3, \oy) -- (\dx+\ox+0.3, \dy+\oy)
    node[midway, right, font=\small\bfseries] {Batch $B$};

  % Grid lines on front face (time steps)
  \foreach \y in {0.417,0.833,...,2.083}{
    \draw[gray!40, thin] (0,\y*\dy/2.5) -- (\dx,\y*\dy/2.5);
  }
  % Grid lines (features)
  \foreach \x in {0.393,0.786,...,5.107}{
    \draw[gray!40, thin] (\x,0) -- (\x,\dy);
  }

  % Labels inside cube
  \node[font=\small\bfseries, text=bkblue] at (\dx/2, \dy/2)
    {\texttt{[B, 12, 14]}};

  % Annotation arrows
  \node[right=0.3cm of {(\dx+\ox, \dy/2+\oy/2)}, font=\tiny, text=codegray, text width=2.5cm, align=left]
    {$B = 32$ hoặc $64$\\tùy VRAM GPU};
  \node[below=0.2cm of {(\dx/2,-0.5)}, font=\tiny, text=codegray, text width=5cm, align=center]
    {14 features: 3 GT + 5 KT + 3 KK + 3 PT};
  \node[left=0.3cm of {(-0.4,\dy/2)}, font=\tiny, text=codegray, text width=2.2cm, align=right]
    {12 bước $\times$ 5 phút\\$= 60$ phút lịch sử};
\end{tikzpicture}
\caption{Cấu trúc 3D Tensor đầu vào của Tầng 2}
\label{fig:tensor_3d}
\end{figure}

Tensor 3D $\mathbf{X} \in \mathbb{R}^{B \times 12 \times 14}$ được xây dựng bằng cơ chế \textbf{sliding window}:

\begin{lstlisting}[caption={Xây dựng 3D Tensor với Sliding Window (PyTorch)},label={lst:tensor_build}]
import torch
import numpy as np
from sklearn.preprocessing import MinMaxScaler

class TrafficTensorBuilder:
    """Xay dung 3D Tensor [Batch, TimeSteps, Features] cho STGCN."""

    def __init__(self, time_steps: int = 12, features: int = 14):
        self.T = time_steps   # 12 buoc x 5 phut = 60 phut lich su
        self.F = features     # 14 bien dac trung
        self.scaler = MinMaxScaler(feature_range=(0, 1))
        self._scaler_fitted = False

    def fit_scaler(self, raw_data: np.ndarray) -> None:
        """Fit scaler CHI TREN tap training. KHONG goi lai khi inference."""
        if raw_data.shape[1] != self.F:
            raise ValueError(f"Expected {self.F} features, got {raw_data.shape[1]}")
        self.scaler.fit(raw_data)
        self._scaler_fitted = True

    def build(self, raw_window: np.ndarray) -> torch.Tensor:
        """
        Tham so:
            raw_window: [N_timesteps, N_features] - du lieu thu trong cua mot nut
        Tra ve:
            tensor: [1, 12, 14] - Tensor san sang cho model
        """
        if not self._scaler_fitted:
            raise RuntimeError("Scaler chua duoc fit. Goi fit_scaler() truoc.")

        # 1. Lay 12 buoc gan nhat
        window = raw_window[-self.T:, :]           # [12, 14]

        # 2. Chuan hoa MinMax
        normalized = self.scaler.transform(window) # [12, 14]

        # 3. Chuyen sang Tensor va them batch dim
        tensor = torch.tensor(normalized, dtype=torch.float32)
        return tensor.unsqueeze(0)                 # [1, 12, 14]

    def build_batch(self, windows: list) -> torch.Tensor:
        """Xay dung batch tensor cho nhieu nut giao."""
        tensors = [self.build(w) for w in windows]
        batch = torch.cat(tensors, dim=0)          # [B, 12, 14]
        assert batch.shape[1:] == (12, 14), f"Sai shape: {batch.shape}"
        return batch
\end{lstlisting}

\section{Kiến trúc Lưu trữ và Đồng bộ Dữ liệu}

Để đảm bảo đồng bộ dữ liệu từ nhiều nguồn (CCTV và cảm biến có thể có độ trễ khác nhau), pipeline sử dụng kiến trúc \textbf{Event-Driven + Time-Window Buffering}:

\begin{enumerate}[leftmargin=1.5cm, itemsep=4pt]
  \item Tất cả messages được publish lên message broker (MQTT) với timestamp chính xác.
  \item Time-window buffer (5 phút) tích lũy dữ liệu từ tất cả sensors tại một nút.
  \item Sau khi cửa sổ đóng, data aggregator tính toán trung bình và ghép thành vector 14 chiều.
  \item Nếu một cảm biến thiếu dữ liệu, áp dụng chiến lược imputation tương ứng.
  \item Vector hoàn chỉnh được ghi vào InfluxDB với key `\texttt{node\_id + timestamp}'.
\end{enumerate}
